% Two sufficient routes each combine generalization control and empirical optimality.
% Volume three architecture figures. Independent export only; not book fonts.
\colorlet{MutedBlueFill}{FigureBlueFill}
\colorlet{MutedGreenFill}{FigureGreenFill}
\colorlet{MutedPurple}{FigurePurple}
\colorlet{MutedPurpleFill}{FigurePurpleFill}
\tikzset{
 v3 picture/.style={x=1mm,y=1mm,font=\figurefont\fontsize{8}{10}\selectfont,text=FigureInk,line width=.65pt},
 v3 box/.style={draw=FigureBlue,fill=white,rounded corners=1.8pt,line width=.7pt,inner sep=2.5pt,align=center},
 v3 data/.style={v3 box,fill=FigureBlueFill},
 v3 model/.style={v3 box,draw=FigureGreen,fill=FigureGreenFill},
 v3 output/.style={v3 box,draw=FigureYellow,fill=FigureYellowFill},
 v3 risk/.style={v3 box,draw=FigureRed,fill=FigureRedFill},
 v3 group/.style={draw=FigureBlue,rounded corners=1.8pt,line width=.65pt,inner sep=3mm},
 v3 flow/.style={draw=FigureStroke,line width=.75pt,-{Stealth[length=1.65mm,width=1.05mm]},rounded corners=1pt},
 v3 guide/.style={draw=FigureMuted,line width=.6pt,dashed},
 v3 feedback/.style={v3 flow,draw=FigureGreen,dashed},
 v3 title/.style={draw=none,fill=none,font=\figurefont\fontsize{8}{10}\selectfont,align=center,inner sep=1pt},
 v3 note/.style={draw=none,fill=none,font=\FigureNoteFont\fontsize{7.5}{9.5}\selectfont,text=FigureMuted,align=center,inner sep=1pt},
 v3 label/.style={draw=none,fill=none,font=\figurefont\fontsize{8}{10}\selectfont,align=center,inner sep=1pt},
 v3 wire/.style={draw=FigureMuted,line width=.45pt},
 v3 point/.style={circle,draw=FigureBlue,fill=FigureBlueFill,minimum size=2.3mm,inner sep=0pt,line width=.55pt}
}
% Bring pre-existing native TikZ geometry to the same typography and scale.
\tikzset{
 foundation picture/.append style={font=\figurefont\fontsize{8}{10}\selectfont,line width=.65pt},
 foundation node/.append style={rounded corners=1.8pt,line width=.7pt,font=\figurefont\fontsize{8}{10}\selectfont},
 foundation heading/.append style={draw=none,fill=none,font=\figurefont\fontsize{8}{10}\selectfont},
 foundation note/.append style={draw=none,fill=none,font=\FigureNoteFont\fontsize{7.5}{9.5}\selectfont},
 foundation flow/.append style={line width=.75pt,-{Stealth[length=1.65mm,width=1.05mm]}},
 foundation frame/.append style={rounded corners=1.8pt,line width=.7pt},
 foundation edge label/.append style={draw=none,fill=none,font=\figurefont\fontsize{8}{10}\selectfont}
}
% Official Lucide outline icons, consistently used across diagrams. Size in mm.
\newcommand{\VIcon}[5]{\node[inner sep=0pt] at (#2,#3) {\includegraphics[width=#4mm]{figures/03-foundations/00-shared/tikz/icons/#1-#5.pdf}};}
% A neutral network glyph (not a claim about an actual model architecture).
% Coordinates denote its centre; scale 1 gives 13 x 8 mm.
\newcommand{\VNet}[4]{\begin{scope}[shift={(#1,#2)},scale=#3]
 \foreach \a in {-3,0,3}{\foreach \b in {-3,0,3}{\draw[v3 wire] (-5,\a)--(0,\b);}}
 \foreach \a in {-3,0,3}{\draw[v3 wire] (0,\a)--(5,0);}
 \foreach \a in {-3,0,3}{\draw[draw=#4,fill=#4,line width=.55pt] (-5,\a) circle(.9);\draw[draw=#4,fill=white,line width=.55pt] (0,\a) circle(.9);}
 \draw[draw=#4,fill=#4,line width=.55pt] (5,0) circle(.9);
\end{scope}}
\newcommand{\VDots}[4]{\foreach \k in {0,1,2,3}{\fill[#4] ({#1+2.7*\k},#2) circle(#3);}}
% Feature vector/tensor elements are scientific glyphs, not decorative icons.
\newcommand{\VVector}[4]{\foreach \k in {0,...,4}{\fill[#4,rounded corners=.25pt] ({#1+2.2*\k},#2) rectangle ++(1.5,#3);}}

\begin{tikzpicture}[v3 picture]
\path[use as bounding box] (0,0) rectangle (163,76);
\node[v3 title,text=FigureBlue] at (40,72) {空间层面的控制};
\node[v3 title,text=FigureGreen] at (123,72) {算法层面的控制};
\foreach \x/\c in {3/FigureBlue,42/FigureBlue,86/FigureGreen,125/FigureGreen}{\draw[v3 group,draw=\c] (\x,39) rectangle ++(35,27);}
\node[v3 label] at (20.5,62) {整个空间一致收敛};
\node[v3 label] at (59.5,62) {经验风险接近最优};
\node[v3 label] at (103.5,62) {稳定性代价 $\beta_n$};
\node[v3 label] at (142.5,62) {经验比较误差 $\rho_n$};
% Qualitative hypothesis curves, with no numerical axes or implied measurements.
\draw[v3 flow] (6,43)--(34,43);\draw[v3 flow] (6,43)--(6,57);
\draw[FigureBlue,line width=.7pt] (6,47) .. controls (13,42) and (18,49) .. (23,52) .. controls (28,56) and (31,55) .. (34,52);
\draw[FigureBlue,dashed] (6,50) .. controls (13,46) and (17,57) .. (24,50) .. controls (28,45) and (31,46) .. (34,48);
\draw[FigurePurple,dotted] (6,53) .. controls (10,61) and (15,42) .. (22,47) .. controls (26,52) and (31,50) .. (34,47);
\foreach \off in {0,83}{
 \draw[v3 flow] ({45+\off},43)--({74+\off},43);\draw[v3 flow] ({45+\off},43)--({45+\off},57);
 \draw[FigureBlue,line width=.7pt] ({45+\off},46) .. controls ({54+\off},48) and ({60+\off},56) .. ({73+\off},52);
 \foreach \x/\y in {48/47,50/49,52/45,55/51,57/47,60/54,62/51,65/56,68/54,70/50,73/51}{\fill[FigureBlue] ({\x+\off},\y) circle(.55);}}
% Adjacent sample rows, one replacement, and algorithm outputs.
\foreach \y in {47,54}{\foreach \j in {0,...,4}{\fill[FigureBlue,rounded corners=.25pt] ({88+3.1*\j},\y) rectangle ++(2.1,2.4);}}
\fill[FigureYellow] (94.2,54) rectangle ++(2.1,2.4);\fill[FigureGreen] (94.2,47) rectangle ++(2.1,2.4);
\VIcon{settings}{111}{52}{6}{ink}
\draw[v3 flow] (104,55)--(107,53);\draw[v3 flow] (104,48)--(107,50);
\draw[v3 flow] (114,53)--(116,55);\draw[v3 flow] (114,50)--(116,48);
\fill[FigureBlue] (118,55) circle(1.1);\fill[FigureGreen] (118,48) circle(1.1);
% Only conjunction produces the final guarantee.
\foreach \x/\c in {40/FigureBlue,123/FigureGreen}{
 \draw[v3 wire,draw=\c] ({\x-19.5},39)--({\x-19.5},34)--({\x+19.5},34)--({\x+19.5},39);
 \draw[v3 flow,draw=\c] (\x,34)--(\x,29);
 \node[v3 box,draw=\c,minimum width=35mm,minimum height=10mm] (combine\x) at (\x,24) {两项前提结合};
 \draw[v3 flow,draw=\c] (\x,19)--(\x,13);
}
\node[v3 box,minimum width=74mm,minimum height=12mm] at (40,7) {期望风险接近最优\\{\FigureNoteFont\fontsize{7.5}{9.5}\selectfont 需要同时满足两个前提}};
\node[v3 box,draw=FigureGreen,minimum width=74mm,minimum height=12mm] at (123,7) {平均超额风险 $\leq\beta_n+\rho_n$};
\end{tikzpicture}
